% Source presentations for register row N-c-003.

\subsection{Locality in Libkin and the finite Hanf provenance boundary}
\label{sec:p09-presentation-n-c-003}

Libkin's presentation keeps three notions apart.  First, a query is
\emph{Hanf-local} when a bijection matching all pointed neighbourhoods up to
some fixed radius preserves its answers.  Second, a non-Boolean query is
\emph{Gaifman-local} when the answer on a tuple is determined by the induced
neighbourhood of that tuple at some fixed radius.  Third, Gaifman normal form
is a syntactic representation of formulas by local formulas and separated-
witness sentences.  Global comparison of all pointed neighbourhoods, local
dependence around one answer tuple, and a formula normal form are therefore
different assertions \cite[defs.~4.6--4.7 and thm.~4.22]{Libkin2004}.
When the first or second radius exists, write \(\operatorname{hlr}(Q)\) or
\(\operatorname{lr}(Q)\) for its least value, respectively.  These are
numerical abbreviations for the two least radii, introduced only because the
source compares them quantitatively below.

For the first two notions, Libkin supplies quantitative game arguments.  If a
non-Boolean query \(Q\) is Hanf-local, then it is Gaifman-local and
\[
 \operatorname{lr}(Q)\leq 3\operatorname{hlr}(Q)+1.
\]
If \(Q\) is defined by a first-order formula of quantifier rank at most \(k\),
then
\[
 \operatorname{hlr}(Q)\leq \frac{3^k-1}{2},
 \qquad
 \operatorname{lr}(Q)\leq \frac{3^{k+1}-1}{2}.
\]
The proofs maintain partial isomorphisms between successively shrinking tuple
neighbourhoods in the Ehrenfeucht--Fra\"iss\'e game
\cite[thms.~4.11--4.12]{Libkin2004}.  These semantic locality theorems do not by
themselves produce the normal-form syntax.

The normal-form statement is separate.  For a relational vocabulary, every
first-order formula \(\varphi(\bar x)\) is equivalent to a Boolean combination
of formulas local around \(\bar x\) and sentences of the form
\[
 \exists x_1\cdots\exists x_s
 \left(
   \bigwedge_{i=1}^{s}\alpha^{(r)}(x_i)
   \land
   \bigwedge_{1\leq i<j\leq s}d(x_i,x_j)>2r
 \right),
\]
where \(\alpha^{(r)}\) is local around its displayed variable.  Sentences need
only the separated-witness components.  For a formula of quantifier rank
\(k\) with \(n\) free variables, Libkin records the effective bounds
\(r\leq7^k\) and \(s\leq k+n\).  The source states this as Theorem~4.22 but
assigns its proof to Exercise~4.9; the proved semantic locality theorems above
must not be relabelled as a proof of that exercise
\cite[pp.~60--63]{Libkin2004}.  The historical source is
\textcite{Gaifman1982}; the exact syntax, bounds, and exercise status just used
are those of Libkin's later presentation.

For finite structures of bounded degree, Definition~4.23 replaces a bijection
of neighbourhood occurrences by threshold comparison: each rooted radius type
has the same number of realizations below a threshold, or at least the threshold
in both structures.  Theorem~4.24 proves that for every quantifier rank \(k\)
and degree bound \(\ell\) there are a radius \(d\) and threshold \(m\) such
that threshold-equivalent structures of degree at most \(\ell\) are
\(k\)-round EF-equivalent.  Its proof takes
\(r_0=0\), \(r_{i+1}=3r_i+1\), and \(d=r_{k-1}\), then chooses \(m\) large
enough for the finitely many bounded-degree neighbourhood types and preserves
isomorphic shrinking tuple-neighbourhoods during play
\cite[pp.~61--62]{Libkin2004}.

The finite threshold theorem has a later primary provenance distinct from
Hanf's 1965 chapter.  Fagin, Stockmeyer, and Vardi define
\((d,m)\)-equivalence by the same below-threshold/equally-large comparison and
prove that, for \(r,f>0\), finite structures of degree at most \(f\) which are
\((d,m)\)-equivalent are \(r\)-round EF-equivalent; their proof may take
\(d=3^{r-1}\) and \(m=r f^{d-1}\)
\cite[thm.~4.3]{FaginStockmeyerVardi1995}.  They then describe Hanf's original
result as an infinite, locally finite elementary-equivalence criterion using
exact finite equality of every neighbourhood count or infinitude in both
structures, and expressly warn that the two results do not directly imply one
another \cite[lem.~4.5]{FaginStockmeyerVardi1995}.  Thus the historical name
``Hanf locality'' does not authorize attribution of the finite threshold
theorem to \textcite{Hanf1965}.

The cycle constructions in Libkin show the use of these distinctions.  Two
long cycles can agree on all bounded neighbourhoods while only their disjoint
union pattern changes connectivity; loop-marked cycles can agree only after
thresholding while differing in a global counting property.  The first kind of
example separates a global property from bijection-based locality, and the
second uses bounded degree and threshold counts
\cite[pp.~48 and 62--63]{Libkin2004}.  Neither example says that an arbitrary
first-order formula is itself one local formula.

\paragraph{Original-source boundaries.}
The registered Hanf 1965 chapter and Gaifman 1982 article remain separate
presentations.  Their bibliographic records establish identity, and Gaifman's
institutional abstract supports the qualitative local-normal-form description,
but the source pages needed to reconstruct their own notation, theorem numbers,
proofs, and examples were not available for verification.  No quantitative
bound or later finite threshold convention is therefore retrojected into those
originals \cite{Hanf1965,Gaifman1982}.
